\chapter{第十一届大唐杯本科B组省赛第一场}
\fancyhead[L]{\color{H6}\kaishu\faIcon{atom}\;第十一届大唐杯本科B组省赛第一场}
\fancyhead[R]{\color{H6}\kaishu\rightmark\,}

\date{2024年4月20日}{}{\href{https://github.com/xubenshan/dtcup}{\textbf{才疏学浅的小熊}}}
{}
{https://www.xiaohongshu.com/user/profile/6192219d0000000021028647}{小红书}
{https://space.bilibili.com/692546609}{B站账号}




\section{单选题（共40小题，共120分）}



\begin{choice}{C}[]
    对于三层交换机，若要配置物理端口的IP地址，以下配置方法中，不正确的是
    \begin{tasks}(1)
        \task 对于Cisco或锐捷的三层交换机，在接口配置模式下，执行noswitchport命令将端口设置为三层端口，然后使用iPaddress命令设置端口的IP地址即可。
        \task 对于Cisco或锐捷的三层交换机，实现创建一个VLAN，并在该VLAN接口上配置IP地址，然后将要配置IP地址的端口划归属于该VLAN即可。
        \task 对于华为或华三的三层交换机，在接口配置模式下，执行undoswitchport命令将端口设置为三层端口，然后使用iPaddress命令设置端口的IP地址即可
        \task 对于华为或华三的三层交换机，首先创建一个VLAN，并在该VLAN接口上配置IP地址，然后将要配置IP地址的端口划归属于该VLAN即可
    \end{tasks}
\end{choice}


\begin{choice}{D}[]
    基站传输规划表中，30位掩码对应的子网掩码是
    \begin{tasks}(4)
        \task 255.255.255.248
        \task 255.255.255.240
        \task 255.255.255.0
        \task 255.255.255.252
    \end{tasks}
\end{choice}




\begin{choice}{B}[]
    在IPD模式的流程中，哪个阶段开始产品正式立项
    \begin{tasks}(4)
        \task 第二阶段
        \task 第一阶段
        \task 第四阶段
        \task 第三阶段
    \end{tasks}
\end{choice}



\begin{choice}{A}[]
    以下(\quad)产品成本核算方法属于财务会计范畴
    \begin{tasks}(4)
        \task 完全成本法
        \task 作业变动成本法
        \task 作业成本法
        \task 变动成本法
    \end{tasks}
\end{choice}


\begin{choice}{A}[]
    在通信系统数字化移动信道中，交织编码作用为
    \begin{tasks}(4)
        \task 纠正突发差错
        \task 应对打孔误码
        \task 纠正随机差错
        \task 应对冗余误码
    \end{tasks}
\end{choice}

\begin{choice}{A}[]
    5G NR协议栈中，以下选项中属于5G NG口控制面协议栈组成的是
    \begin{tasks}(4)
        \task SCTP
        \task RRC
        \task GTP\_U
        \task SDAP
    \end{tasks}
\end{choice}




\begin{choice}{C}[]
    在5G网络中提到的毫米波的频率范围是
    \begin{tasks}(2)
        \task 300-3000MHZ
        \task 30-300MHZ
        \task 30-300GHZ
        \task 3-30GHZ
    \end{tasks}
\end{choice}


\begin{choice}{A}[]
    以下不属于决策树求解算法的选项是
    \begin{tasks}(4)
        \task 梯度下降法
        \task CART算法
        \task C4.5
        \task ID3
    \end{tasks}
\end{choice}


\begin{choice}{A}[]
    工程经济学中经济的含义主要是指
    \begin{tasks}(4)
        \task 资源的合理利用
        \task 经济制度
        \task 生产关系
        \task 经济基础
    \end{tasks}
\end{choice}


\begin{choice}{C}[]
    预计6G通信将支持五个应用场景，其中安全可靠低时延通信场景可简称为
    \begin{tasks}(2)
        \task BigCom
        \task eMBB-Plus
        \task SURLLC
        \task 3D-InteCom
    \end{tasks}
\end{choice}

\begin{choice}{C}[]
    5G NR中RRC建立过程本质上是建立空口某一承载过程，该承载是
    \begin{tasks}(4)
        \task SRB3
        \task SRB2
        \task SRB1
        \task SRB0
    \end{tasks}
\end{choice}


\begin{choice}{D}[]
    在5G基站安装中，DCPD与电源柜之间的电源线规格应为
    \begin{tasks}(4)
        \task 16方
        \task 8方
        \task 10方
        \task 25方
    \end{tasks}
\end{choice}


\begin{choice}{C}[]
    华为Mate60pro的卫星通话功能采用中国自主研制的卫星移动通信系统由（\quad）运营
    \begin{tasks}(4)
        \task 中国联通
        \task 中国移动
        \task 中国电信
        \task 中国卫通
    \end{tasks}
\end{choice}


\begin{choice}{C}[]
    下列选项中哪个不是5G NR中定义的子载波间隔
    \begin{tasks}(4)
        \task 30KHZ
        \task 15KHZ
        \task 90KHZ
        \task 60KHZ
    \end{tasks}
\end{choice}

\begin{choice}{C}[]
    3GPP发布全球统一标准的车联网通信技术是指
    \begin{tasks}(4)
        \task V2P
        \task V2I
        \task C-V2X
        \task DSRC
    \end{tasks}
\end{choice}

\begin{choice}{A}[]
    5G网络架构中，CU/DU分离场景下，回传一般是指（\quad）之间的数据传送
    \begin{tasks}(2)
        \task CU-5GC
        \task AAU-DU
        \task DU-CU
        \task AAU-BBU
    \end{tasks}
\end{choice}



\begin{choice}{B}[]
    多个GPS天线之间的间距不能过近，否则天线之间可能会产生反射干扰，两个GPS天线间距大于（\quad）米
    \begin{tasks}(4)
        \task 4
        \task 2
        \task 1
        \task 3
    \end{tasks}
\end{choice}

\begin{choice}{C}[]
    无线对讲通话设备主要用于提供实时语音服务，其主要作用不包括
    \begin{tasks}(4)
        \task 应急通信
        \task 团队协作
        \task 抗干扰通信
        \task 便捷通信
    \end{tasks}
\end{choice}


\begin{choice}{C}[]
    4G\;LTE有（\quad）个物理层小区ID
    \begin{tasks}(4)
        \task 512
        \task 1008
        \task 504
        \task 256
    \end{tasks}
\end{choice}



\begin{choice}{C}[]
    5G\;NR系统中，对于SA场景，UE\;ip\;address\;allocation功能由（\quad）功能模块实现
    \begin{tasks}(4)
        \task UPF
        \task MME
        \task SMF
        \task AMF
    \end{tasks}
\end{choice}



\begin{choice}{B}[]
    随机森林是决策树的集成，是一种(\quad)方法
    \begin{tasks}(4)
        \task pasting
        \task Bagging
        \task Boosting
        \task stacking
    \end{tasks}
\end{choice}




\begin{choice}{C}[]
    5G标准的第四版对应为3GPP的(\quad)版本
    \begin{tasks}(4)
        \task R16
        \task R15
        \task R18
        \task R19
    \end{tasks}
\end{choice}

\begin{choice}{A}[]
    5G NR系统中，一个ss/PBCH\;block包含(\quad)个OFDM\;symbols
    \begin{tasks}(4)
        \task 4
        \task 2
        \task 1
        \task 3
    \end{tasks}
\end{choice}

\begin{choice}{D}[]
    5G\;NR系统中，物理小区组ID标识范围为
    \begin{tasks}(4)
        \task 0-2
        \task 0-503
        \task 0-3
        \task 0-335
    \end{tasks}
\end{choice}


\begin{choice}{B}[]
    5G\;NR系统R15版本中，PDSCH支持的最大调制方式为
    \begin{tasks}(4)
        \task 512QAM
        \task 256QAM
        \task 64QAM
        \task 128QAM
    \end{tasks}
\end{choice}


\begin{choice}{D}[]
    在3GPP R16中进一步增强了动态频谱共享技术，允许LTE和NR共享相同的载波，动态频谱共享简称为
    \begin{tasks}(4)
        \task TDD
        \task FDD
        \task CSS
        \task DSS
    \end{tasks}
\end{choice}



\begin{choice}{B}[]
    新产品开发的启动和终结由（\quad）来决定
    \begin{tasks}(4)
        \task IPD
        \task IPMT
        \task PDT
        \task SPT
    \end{tasks}
\end{choice}


\begin{choice}{A}[]
    局域网LAN内的节点必须有网络接口设备和
    \begin{tasks}(2)
        \task 主机地址
        \task 包含所有网络节点的表
        \task 网络账户
        \task 共享资源
    \end{tasks}
\end{choice}


\begin{choice}{C}[]
    以下说法错误的是
    \begin{tasks}(1)
        \task IPD模式对新产品的开发是以客户为中心的市场行为，追求的是客户对新产品的满意度
        \task IPD模式对新产品的开发是以技术创新为目的，追求的是新产品的市场应用
        \task IPD模式采用技术驱动和市场驱动相结合，将新产品开发作为技术型决策，而非投资性决策
        \task IPD模式对新产品的开发是作为投资进行管理的，追求的是投资回报
    \end{tasks}
\end{choice}


\begin{choice}{A}[]
    GBDT和AdaBoost的区别是
    \begin{tasks}(1)
        \task AdaBoost在每个迭代中调整实例权重，而GBDT让新的预测器针对前一个预测器的残差进行拟合
        \task 二者无区别
        \task 预测器不同
        \task GBDT在每个迭代中调整实例权重，而AdaBoost让新的预测器针对前一个预测器的残差进行拟合
    \end{tasks}
\end{choice}


\begin{choice}{A}[]
    移动通信每十年一代，6G技术是5G代际更新的一个新技术，预期其商用时间是
    \begin{tasks}(2)
        \task 2030左右
        \task 2024左右
        \task 2025左右
        \task 2029左右
    \end{tasks}
\end{choice}


\begin{choice}{A}[]
    大唐5G基站安装时，当BBU安装+AAU设备安装2台时，上级输入空开建议为
    \begin{tasks}(2)
        \task 2路80A
        \task 2路100A
        \task 1路80A
        \task 1路63A
    \end{tasks}
\end{choice}



\begin{choice}{B}[]
    在5G NR中，部分系统消息可通过Ondemand方式发送，网络可以避免周期性广播这些消息，以下选项中，不属于该类系统消息的是
    \begin{tasks}(4)
        \task SIB2
        \task SIB1
        \task SIB4
        \task SIB3
    \end{tasks}
\end{choice}


\begin{choice}{B}[]
    5G基站设备中，当DCPD至AAU电源线长度大于50米，小于等于80米，直流电源线应使用
    \begin{tasks}(2)
        \task $2*25mm^2$
        \task $2*16mm^2$
        \task $2*10mm^2$
        \task $2*6mm^2$
    \end{tasks}
\end{choice}


\begin{choice}{B}[]
    以下哪些不是车联网汽车必须具备的条件
    \begin{tasks}(2)
        \task 融合网络技术
        \task 支持环境温湿度检测
        \task 搭配控制器
        \task 搭载传感器
    \end{tasks}
\end{choice}


\begin{choice}{C}[]
    根据不同分类标准通信可对应不同的分类，则按照哪一种分类标准可分为数字通信和模拟通信
    \begin{tasks}(2)
        \task 通信方式
        \task 传输媒质
        \task 信号特征
        \task 通信业务
    \end{tasks}
\end{choice}


\begin{choice}{C}[]
    根据香农公式，属于不能提高信道容量的方法是
    \begin{tasks}(2)
        \task 采用更高效的调制方式
        \task 扩展信道带宽
        \task 降低信噪比
        \task 使用多天线技术
    \end{tasks}
\end{choice}


\begin{choice}{D}[]
    3GPP未来将向6G逐步演进，3GPP的NTN标准是从哪一版本开始启动，一直朝着将卫星纳入3GPP技术规范的目标在前进
    \begin{tasks}(4)
        \task R14
        \task R13
        \task R16
        \task R15
    \end{tasks}
\end{choice}


\begin{choice}{A}[]
    5G场景中，智能抄表业务对哪项指标要求较高
    \begin{tasks}(4)
        \task 连接数
        \task 速率
        \task 时延
        \task 峰值速率
    \end{tasks}
\end{choice}


\begin{choice}{D}[]
    6G网络将在智享生活，智赋生产等方面催生全新的应用场景，这些场景将需要很多全新的网络功能和能力，其中不包括
    \begin{tasks}(2)
        \task 数字孪生
        \task 太比特级的峰值速率
        \task 超过1000km/h的移动速度
        \task 毫秒级的时延体验
    \end{tasks}
\end{choice}



\section{多选题（共30小题，共120分）}

\begin{choice}{\;BD\;}[]
    在基站勘察设计过程中，以下不属于勘察流程的有
    \begin{tasks}(2)
        \task 天馈勘察
        \task 基站设备安装
        \task 勘察前准备
        \task 仔细研读设备说明书
    \end{tasks}
\end{choice}

\begin{choice}{\;ABD\;}[]
    智能网联汽车环境感知传感器属于信息收集单元的重要组成部分，其信息收集能力主要有什么决定
    \begin{tasks}(2)
        \task 传感器的数量
        \task 传感器感知范围
        \task 车辆所在位置
        \task 传感器感知精度
    \end{tasks}
\end{choice}

\begin{choice}{\;ABCD\;}[]
    在未来6G移动通信系统中，AI在物理层主要可用于什么方面
    \begin{tasks}(2)
        \task 端到端链路设计
        \task CSI处理
        \task 接收机设计
        \task 发射机设计
    \end{tasks}
\end{choice}

\begin{choice}{\;CD\;}[]
    市场细分的目的包括
    \begin{tasks}(2)
        \task 分析出潜在的竞争对手
        \task 分析出本产品的替代性技术
        \task 确定产品的客户是什么样的群体
        \task 确定产品的客户购买本企业产品的理由是什么
    \end{tasks}
\end{choice}


\begin{choice}{\;ABCD\;}[]
    以下属于集成学习方法的选项是
    \begin{tasks}(4)
        \task Boosting
        \task stacking
        \task Bagging
        \task Voting
    \end{tasks}
\end{choice}

\begin{choice}{\;BCD\;}[]
    无线对讲通话设备主要用于提供实时语音通信服务，其主要作用包括
    \begin{tasks}(4)
        \task 抗干扰通信
        \task 团队协作
        \task 应急通信
        \task 便捷通信
    \end{tasks}
\end{choice}

\begin{choice}{\;BD\;}[]
    在工业互联网标识解析中，以下哪个说法是正确的
    \begin{tasks}(1)
        \task 工业互联网标识解析仅适用于制造业，不适用于其他行业
        \task 工业互联网标识解析能够为企业提供产品追溯，防伪防窜，供应链管理等服务
        \task 工业互联网标识解析的核心是标识注册，不需要考虑标识的查询和管理
        \task 工业互联网标识解析系统可以支持多种标识编码体系，实现不同系统间的互操作性
    \end{tasks}
\end{choice}


\begin{choice}{\;BCD\;}[]
    5G网络部署Massive\;MIMO主要用于解决哪些问题
    \begin{tasks}(2)
        \task 穿透损耗问题
        \task 大容量问题
        \task 深度覆盖问题
        \task 天面不足问题
    \end{tasks}
\end{choice}


\begin{choice}{\;ABCD\;}[]
    通信领域中，分集的方式包括
    \begin{tasks}(4)
        \task 时间分集
        \task 空间分集
        \task 频率分集
        \task 极化分集
    \end{tasks}
\end{choice}


\begin{choice}{\;ACD\;}[]
    5G\;R15版本，信道状态信息CSI可能包括以下选项中的哪些内容
    \begin{tasks}(4)
        \task CQI
        \task SINR
        \task RI
        \task PMI
    \end{tasks}
\end{choice}


\begin{choice}{\;ACD\;}[]
    以下关于经济学的说法正确的是
    \begin{tasks}(2)
        \task 解决人的欲望与现实资源之间的矛盾
        \task 用来指导人们赚钱
        \task 由日本学者神田孝平翻译而来
        \task 其中也包括治理国家的含义
    \end{tasks}
\end{choice}

\begin{choice}{\;ABD\;}[]
    关于GPS的安装规范，以下说法正确的有
    \begin{tasks}(1)
        \task 放大器安装在离GPS天线50m-150m处，安装在墙上或走线架上
        \task 一定要在避雷针的45度防雷保护范围内
        \task GPS安装不需要接地
        \task GPS天线宜远离其他发射或接收设备，不要安装在微波天线，高压线下方，避免发射天线的辐射方向对准GPS天线
    \end{tasks}
\end{choice}

\begin{choice}{\;ABCD\;}[]
    工业互联网典型应用包括
    \begin{tasks}(2)
        \task 面向社会化生产的资源优化配置与协同
        \task 面向产品全生命周期的管理与服务优化
        \task 面向工业现场的生产过程优化
        \task 面向企业运营的管理决策优化
    \end{tasks}
\end{choice}

\begin{choice}{\;ABCD\;}[]
    在5G\;NR中搜索空间的类型有多种，各个搜索空间支持的RNTI类型不尽相同，其中type2-PDCCH\;CSS支持的RNTI类型可能为
    \begin{tasks}(4)
        \task UCI
        \task CORESET
        \task Search\;Space
        \task DCI
    \end{tasks}
\end{choice}

\begin{choice}{\;ABCD\;}[]
    车联网中，环境感知主要包括
    \begin{tasks}(1)
        \task 行人感知，主要判断车辆行驶前方是否有行人，包括白天行人识别，夜晚行人识别，被障碍物遮挡的行人识别
        \task 车辆本身状态感知、包括行驶速度、行驶方向、行驶状态、车辆位置等
        \task  道路感知，包括道路类型检测、道路标线识别、道路状况判断、是否偏离行驶轨迹等
        \task 交通信号感知、自动识别交叉路口的信号灯、如何高效通过交叉路口等
    \end{tasks}
\end{choice}

\begin{choice}{\;ABCD\;}[]
    在5G中终端通过解读SIB1消息可获取很多信息，以下属于SIB1携带内容
    \begin{tasks}(2)
        \task 服务小区公共配置
        \task 小区选择信息
        \task SI调度信息
        \task 小区接入相关信息
    \end{tasks}
\end{choice}

\begin{choice}{\;ABC\;}[]
    下列选项中属于5G参考信号的是
    \begin{tasks}(4)
        \task SRS
        \task DMRS
        \task PTRS
        \task CRS
    \end{tasks}
\end{choice}

\begin{choice}{\;BCD\;}[]
    5G\;NR中包括哪些Raster栅格
    \begin{tasks}(2)
        \task measure Raster
        \task Global Raster
        \task Channel Raster
        \task SS Raster
    \end{tasks}
\end{choice}

\begin{choice}{\;ABCD\;}[]
    在5G中SSB\;Case\;C类型中，频段大于2.4G时，SSB在一个无线帧中数量为8，则对应8个SSB分别存在于哪些时隙上
    \begin{tasks}(4)
        \task 时隙1
        \task 时隙3
        \task 时隙2
        \task 时隙0
    \end{tasks}
\end{choice}

\begin{choice}{\;AB\;}[]
    自由空间的传播模型的无线电波损耗与什么有关
    \begin{tasks}(2)
        \task 传播距离
        \task 电波频率
        \task 一般城区、郊区
        \task 遮挡
    \end{tasks}
\end{choice}

\begin{choice}{\;ABC\;}[]
    以下选项中哪些属于5G中切换相关的流程
    \begin{tasks}(2)
        \task 测量过程
        \task 判决过程
        \task 切换执行过程
        \task 小区搜索过程
    \end{tasks}
\end{choice}

\begin{choice}{\;BCD\;}[]
    以下说法正确的是
    \begin{tasks}(1)
        \task 只要是自己的技术，自己的筹资，自己的开发人员，开发出来的新产品就一定具有自主知识产权
        \task 发明专利的申请要符合新颖性的要求，意即早于他人知晓，提出创新成果，并提交专利申请
        \task 现在用标准必要专利来进行企业专利布局，对竞争对手形成反制，已经成为企业知识产权高端战略的重要组成部分
        \task 若产品设计开发的解决方案中存在他人的授权专利，可以通过向专利权人缴纳专利使用费的方式，化解侵权纠纷
    \end{tasks}
\end{choice}

\begin{choice}{\;BCD\;}[]
    以下属于决策树求解算法的选项是
    \begin{tasks}(4)
        \task 梯度下降法
        \task C4.5
        \task ID3
        \task CART算法
    \end{tasks}
\end{choice}


\begin{choice}{\;ABD\;}[]
    在5G网络中，AMF通过N2接口连接基站，以下属于N2接口协议栈的为
    \begin{tasks}(4)
        \task SCTP
        \task GTP-U
        \task SDAP
        \task NGAP
    \end{tasks}
\end{choice}


\begin{choice}{\;ABCD\;}[]
    在5G到6G，超大规模天线技术持续升级，其中未来重大研究方向的是
    \begin{tasks}(4)
        \task 精准定位
        \task 智能超表面RIS
        \task 太赫兹
        \task 电磁波角动能
    \end{tasks}
\end{choice}


\begin{choice}{\;BD\;}[]
    以下关于6G网络性能指标初步预测情况正确的是
    \begin{tasks}(2)
        \task 用户体验速率100Mbps
        \task UP时延小于0.5ms
        \task 定位精度10m
        \task 峰值速率1Tbps
    \end{tasks}
\end{choice}


\begin{choice}{\;ABCD\;}[]
    6G技术的发展基于信息普惠、绿色能耗、虚拟融合三大发展理念，以下属于6G前沿技术分类的有
    \begin{tasks}(2)
        \task 太赫兹和可见光通信
        \task 天线人工智能
        \task 全息无线电
        \task 智能超表面
    \end{tasks}
\end{choice}


\begin{choice}{\;BD\;}[]
    以下不属于作业成本法优点的是
    \begin{tasks}(2)
        \task 可以获得更准确的产品和产品线成本
        \task 利于管理控制
        \task 有助于改进成本控制
        \task 开发和维护费用低
    \end{tasks}
\end{choice}


\begin{choice}{\;ABCD\;}[]
    AI大数据模型在移动通信网络中的多重作用中包括
    \begin{tasks}(2)
        \task 网络故障诊断与预测
        \task 网络优化
        \task 用户行为分析
        \task 安全保障
    \end{tasks}
\end{choice}


\begin{choice}{\;BCD\;}[]
    5G标准可以分为两个阶段，如下3GPP版本对应属于第二阶段的为
    \begin{tasks}(4)
        \task R17
        \task R18
        \task R20
        \task R19
    \end{tasks}
\end{choice}


\section{判断题（共30小题，共60分）}

\begin{choice}{\;正确\;}[]
    数字对象架构技术是工业互联网标识解析体系搭建的关键支撑技术。
\end{choice}


\begin{choice}{\;正确\;}[]
    6G网络是一种概念性的无线网络移动通信技术，也叫做第六代移动通信技术。
\end{choice}

\begin{choice}{\;错误\;}[]
    由于高频段的波长小，高频的衍射和绕射能力都强于低频。
\end{choice}

\begin{choice}{\;正确\;}[]
    在密集琥珀区，丘陵城市，及有玻璃幕墙建筑的环境中，选址时要注意信号的反射和衍射的影响。
\end{choice}

\begin{choice}{\;错误\;}[]
    5G解决决策问题，AI解决通信问题。
\end{choice}

\begin{choice}{\;错误\;}[]
    智能超表面技术预计能够为6g通信带来显著的容量提升。

\end{choice}


\begin{choice}{\;正确\;}[]
    优化产品成本是经济决策的核心。

\end{choice}

\begin{choice}{\;正确\;}[]
    馈线拐弯应圆滑均匀，弯曲点应尽可能少。

\end{choice}

\begin{choice}{\;正确\;}[]

    5G组网方式可分为独立部署和非独立部署两种。
\end{choice}


\begin{choice}{\;正确\;}[]
    移动通信系统中，在下行链路上，扰码是用于终端用户区分不同的小区。

\end{choice}

\begin{choice}{\;正确\;}[]
    车联网场景中紧急防碰撞预警系统应用属于高频率低时延业务。

\end{choice}


\begin{choice}{\;错误\;}[]
    5G NR中不同系统消息均通过广播方式按照一定周期发送。

\end{choice}


\begin{choice}{\;正确\;}[]
    5G NR中UE处于RRC CONNECT状态时，既可以测量SSB也可以测量CSI信号的RSRP/RSRQ/SINR作为质量评估。

\end{choice}


\begin{choice}{\;错误\;}[]
    5G的基站功能重构为CU和DU两个功能实体，CU主要处理物理层功能和实时性需求的层2功能。

\end{choice}


\begin{choice}{\;错误\;}[]
    电信网按照服务对象分类，可以分为业务网，传输网和支撑网。

\end{choice}


\begin{choice}{\;错误\;}[]
    降低产品变动成本是财务部门的责任。

\end{choice}

\begin{choice}{\;错误\;}[]
    最先进的技术指标就是用户最大的需求。

\end{choice}


\begin{choice}{\;正确\;}[]
    5G NR中随机接入可以分为竞争和非竞争两种。

\end{choice}


\begin{choice}{\;错误\;}[]
    在5G\;NR中pointA的频点是基于PointA的子载波0的中心频率。

\end{choice}


\begin{choice}{\;正确\;}[]
    在Device net现场总线中，实例ID值0表示类本身。

\end{choice}


\begin{choice}{\;正确\;}[]
    无线站址勘测是，需要注意覆盖范围内是否有小于菲涅尔半径的阻挡物。

\end{choice}


\begin{choice}{\;正确\;}[]
    5G NR中初始随机接入失败，终端可以提高发射功率，继续发起随机接入过程。

\end{choice}


\begin{choice}{\;正确\;}[]
    5G NR中UE可以通过SIB3获取同频小区重选的相关消息。

\end{choice}


\begin{choice}{\;正确\;}[]
    无线通信中三种常见效应为阴影效应，远近效应，多普勒效应。其中阴影效应可导致慢衰落，多普勒效应可导致快衰落。

\end{choice}

\begin{choice}{\;错误\;}[]
    5G NR中信令承载DRB被定义为用于传递RRC和NAS消息的无线承载。

\end{choice}

\begin{choice}{\;正确\;}[]
    人工智能AI在5G+核心网方面可以提供网络的管理和优化。

\end{choice}

\begin{choice}{\;错误\;}[]
    C-V2X是融合蜂窝通信与直连通信的车联网通信技术，也叫做NR-V2X。

\end{choice}

\begin{choice}{\;错误\;}[]
    基尼指数代表不确定性，不确定越多，基尼指数越小。

\end{choice}

\begin{choice}{\;错误\;}[]
    对讲机在接收时，会将语音信号转化成电信号。

\end{choice}

\begin{choice}{\;错误\;}[]
    基站开通过程中，网络规划就是板卡规划。

\end{choice}






